\documentclass[11pt]{article}
\usepackage[utf8]{inputenc}
\usepackage[T1]{fontenc}
\usepackage[spanish]{babel}
\usepackage{amsmath, amssymb, amsthm}
\usepackage{braket}
\usepackage[margin=2.5cm]{geometry}
\usepackage{hyperref}

\allowdisplaybreaks
\emergencystretch=2em

\newcommand{\ii}{\mathrm{i}}
\newcommand{\ee}{\mathrm{e}}
\newcommand{\dd}{\mathrm{d}}
\newcommand{\vb}[1]{\mathbf{#1}}
\newcommand{\Ezero}{\mathcal E_0}
\newcommand{\adag}{\hat a^\dagger}
\newcommand{\splus}{\hat\sigma^+}
\newcommand{\sminus}{\hat\sigma^-}

\title{Interacci\'on \'Atomo--Campo}
\author{Ricardo Del Rosal Gardu\~no \\ \small Doctorado en F\'isica, BUAP}
\date{Notas a partir de Schleich, \emph{Quantum Optics in Phase Space}}

\begin{document}
\maketitle

\section{C\'omo construir la interacci\'on}

El objetivo es construir el hamiltoniano que describe la interacci\'on entre un
\'atomo y el campo electromagn\'etico. La idea f\'isica es sencilla: el \'atomo
contiene cargas, y las cargas responden al campo electromagn\'etico. La parte
delicada es escribir esa respuesta de manera compatible con la mec\'anica
cu\'antica, con las simetr\'ias de gauge del electromagnetismo y, m\'as
adelante, con la cuantizaci\'on del propio campo.

En las notas manuscritas se parte de una part\'icula cargada, que puede pensarse
como el electr\'on del \'atomo, de masa \(m_e\), carga \(e\), posici\'on
\(\vb r_e\), y acoplada al campo por medio del potencial vectorial
\(\vb A(\vb r_e,t)\). Se trabaja en el gauge de Coulomb, donde
\(\nabla\cdot\vb A=0\). El campo electromagn\'etico puede describirse en
t\'erminos de los potenciales
\begin{equation}
  \vb B=\nabla\times\vb A,
  \qquad
  \vb E=-\nabla\phi-\frac{\partial\vb A}{\partial t}.
\end{equation}
Por tanto, construir la interacci\'on equivale a entender c\'omo deben entrar
\(\phi\) y \(\vb A\) en la din\'amica de una carga.

Hay dos formas equivalentes de llegar al acoplamiento que se usar\'a en \'optica
cu\'antica. La primera parte de la fuerza de Lorentz y lleva al acoplamiento
m\'inimo, o forma \(\vb A\cdot\vb p\). La segunda usa directamente el momento
dipolar el\'ectrico del \'atomo y lleva a la forma \(-\vb d\cdot\vb E\). Ambas
formas son equivalentes dentro de la aproximaci\'on dipolar y con el centro de
masa fijo, aunque sus elementos de matriz no son iguales t\'ermino a t\'ermino
si no se transforman tambi\'en los estados.

\section{Lagrangiano y fuerza de Lorentz}

El punto de partida experimental es la fuerza de Lorentz sobre una carga \(q\):
\begin{equation}
  \vb F
  =
  q\left(\vb E+\vb v\times\vb B\right).
\end{equation}
Se desea construir un lagrangiano que reproduzca esta fuerza por medio de las
ecuaciones de Euler--Lagrange. Para una carga \(e\), el lagrangiano adecuado es
\begin{equation}
  L
  =
  \frac{1}{2}m\dot{\vb r}^{\,2}
  +e\,\dot{\vb r}\cdot\vb A(\vb r,t)
  -e\,\phi(\vb r,t).
\end{equation}
El primer t\'ermino es la energ\'ia cin\'etica usual. El segundo acopla la
velocidad de la carga con el potencial vectorial. El tercero es el acoplamiento
con el potencial escalar.

Las ecuaciones de Euler--Lagrange, componente por componente, son
\begin{equation}
  \frac{\dd}{\dd t}\left(\frac{\partial L}{\partial\dot r_i}\right)
  -
  \frac{\partial L}{\partial r_i}
  =
  0.
\end{equation}
Primero se calcula la derivada respecto a la velocidad:
\begin{equation}
  \frac{\partial L}{\partial\dot r_i}
  =
  m\dot r_i+eA_i(\vb r,t).
\end{equation}
Al derivar respecto al tiempo debe recordarse que \(A_i\) depende de \(t\) y de
la posici\'on \(\vb r(t)\). Por regla de la cadena,
\begin{equation}
  \frac{\dd}{\dd t}
  \left(\frac{\partial L}{\partial\dot r_i}\right)
  =
  m\ddot r_i
  +e\frac{\partial A_i}{\partial t}
  +e\dot r_j\frac{\partial A_i}{\partial r_j},
\end{equation}
donde se suma sobre el \'indice repetido \(j\). La derivada del lagrangiano
respecto a \(r_i\) es
\begin{equation}
  \frac{\partial L}{\partial r_i}
  =
  e\dot r_j\frac{\partial A_j}{\partial r_i}
  -e\frac{\partial\phi}{\partial r_i}.
\end{equation}
Sustituyendo en Euler--Lagrange,
\begin{align}
  0
  &=
  m\ddot r_i
  +e\frac{\partial A_i}{\partial t}
  +e\dot r_j\frac{\partial A_i}{\partial r_j}
  -e\dot r_j\frac{\partial A_j}{\partial r_i}
  +e\frac{\partial\phi}{\partial r_i},
\end{align}
y, despejando la aceleraci\'on,
\begin{equation}
  m\ddot r_i
  =
  -e\frac{\partial A_i}{\partial t}
  -e\frac{\partial\phi}{\partial r_i}
  -e\dot r_j
  \left(
    \frac{\partial A_i}{\partial r_j}
    -
    \frac{\partial A_j}{\partial r_i}
  \right).
\end{equation}
Los dos primeros t\'erminos son \(eE_i\), porque
\begin{equation}
  E_i
  =
  -\frac{\partial\phi}{\partial r_i}
  -
  \frac{\partial A_i}{\partial t}.
\end{equation}
Queda por identificar el t\'ermino magn\'etico. Usando
\(\vb B=\nabla\times\vb A\),
\begin{align}
  \left(\dot{\vb r}\times\vb B\right)_i
  &=
  \epsilon_{ijk}\dot r_j B_k                                                   \\
  &=
  \epsilon_{ijk}\dot r_j(\nabla\times\vb A)_k                                  \\
  &=
  \epsilon_{ijk}\epsilon_{klm}\dot r_j
  \frac{\partial A_m}{\partial r_l}.
\end{align}
Con la identidad
\begin{equation}
  \epsilon_{ijk}\epsilon_{klm}
  =
  \delta_{il}\delta_{jm}
  -
  \delta_{im}\delta_{jl},
\end{equation}
se obtiene
\begin{align}
  \left(\dot{\vb r}\times\vb B\right)_i
  &=
  \dot r_j
  \left(
    \frac{\partial A_j}{\partial r_i}
    -
    \frac{\partial A_i}{\partial r_j}
  \right)                                                                      \\
  &=
  -\dot r_j
  \left(
    \frac{\partial A_i}{\partial r_j}
    -
    \frac{\partial A_j}{\partial r_i}
  \right).
\end{align}
Por tanto,
\begin{equation}
  m\ddot r_i
  =
  e\left[E_i+\left(\dot{\vb r}\times\vb B\right)_i\right],
\end{equation}
que es justamente la fuerza de Lorentz. Esto confirma que el lagrangiano
anterior contiene el acoplamiento correcto con el campo electromagn\'etico.

\section{Momento can\'onico y acoplamiento m\'inimo}

Una vez construido el lagrangiano, el hamiltoniano se obtiene por transformada
de Legendre. El momento can\'onico conjugado a \(\vb r_e\) es
\begin{equation}
  \vb p_e
  =
  \frac{\partial L}{\partial\dot{\vb r}_e}
  =
  m_e\dot{\vb r}_e+e\vb A(\vb r_e,t).
\end{equation}
De aqu\'i se despeja la velocidad:
\begin{equation}
  \dot{\vb r}_e
  =
  \frac{1}{m_e}
  \left[
    \vb p_e-e\vb A(\vb r_e,t)
  \right].
\end{equation}
Este punto es importante: \(\vb p_e\) no es el momento mec\'anico
\(m_e\dot{\vb r}_e\). El momento mec\'anico es
\begin{equation}
  m_e\dot{\vb r}_e
  =
  \vb p_e-e\vb A(\vb r_e,t).
\end{equation}
Por eso se dice que el campo electromagn\'etico entra mediante el reemplazo
m\'inimo
\begin{equation}
  \vb p_e
  \longrightarrow
  \vb p_e-e\vb A(\vb r_e,t).
\end{equation}

La transformada de Legendre es
\begin{equation}
  H
  =
  \vb p_e\cdot\dot{\vb r}_e-L.
\end{equation}
Sustituyendo la velocidad,
\begin{align}
  H
  &=
  \vb p_e\cdot
  \frac{1}{m_e}\left(\vb p_e-e\vb A\right)
  -
  \frac{1}{2}m_e
  \frac{1}{m_e^2}
  \left(\vb p_e-e\vb A\right)^2
  -
  e
  \frac{1}{m_e}
  \left(\vb p_e-e\vb A\right)\cdot\vb A
  +e\phi .
\end{align}
Para ver las cancelaciones, se agrupan los t\'erminos:
\begin{align}
  H
  &=
  \frac{1}{m_e}\vb p_e^{\,2}
  -\frac{e}{m_e}\vb p_e\cdot\vb A
  -\frac{1}{2m_e}
  \left(
    \vb p_e^{\,2}
    -2e\vb p_e\cdot\vb A
    +e^2\vb A^2
  \right)                                                                     \\
  &\quad
  -\frac{e}{m_e}\vb p_e\cdot\vb A
  +\frac{e^2}{m_e}\vb A^2
  +e\phi .
\end{align}
Los t\'erminos lineales se combinan y queda
\begin{equation}
  H
  =
  \frac{1}{2m_e}
  \left[
    \vb p_e-e\vb A(\vb r_e,t)
  \right]^2
  +e\phi(\vb r_e,t).
\end{equation}
Al desarrollar el cuadrado,
\begin{equation}
  H
  =
  \frac{\vb p_e^{\,2}}{2m_e}
  -
  \frac{e}{m_e}\vb p_e\cdot\vb A(\vb r_e,t)
  +
  \frac{e^2}{2m_e}\vb A^2(\vb r_e,t)
  +
  e\phi(\vb r_e,t).
\end{equation}
Este es el hamiltoniano de acoplamiento m\'inimo.

\section{Gauge de Coulomb y t\'ermino \texorpdfstring{\(\vb A\cdot\vb p\)}{A.p}}

En mec\'anica cu\'antica debe cuidarse el orden de \(\vb p_e\) y
\(\vb A\), porque \(\vb p_e=-\ii\hbar\nabla\) es un operador diferencial. Para
una funci\'on de prueba \(\psi\),
\begin{align}
  \vb p_e\cdot\vb A\,\psi
  &=
  -\ii\hbar\nabla\cdot(\vb A\psi)                                             \\
  &=
  -\ii\hbar(\nabla\cdot\vb A)\psi
  -\ii\hbar\vb A\cdot\nabla\psi .
\end{align}
En el gauge de Coulomb,
\begin{equation}
  \nabla\cdot\vb A=0,
\end{equation}
por lo que
\begin{equation}
  \vb p_e\cdot\vb A\,\psi
  =
  \vb A\cdot\vb p_e\,\psi.
\end{equation}
Dentro de este gauge se puede escribir
\begin{equation}
  \vb p_e\cdot\vb A
  =
  \vb A\cdot\vb p_e.
\end{equation}
El t\'ermino lineal que acopla el campo con el momento de la part\'icula es
\begin{equation}
  \boxed{
  H_{\vb A\cdot\vb p}
  =
  -\frac{e}{m_e}\vb A(\vb r_e,t)\cdot\vb p_e .
  }
\end{equation}
El t\'ermino proporcional a \(\vb A^2\) tambi\'en proviene del acoplamiento
m\'inimo. En campos d\'ebiles suele no ser el protagonista de la absorci\'on y
emisi\'on de un fot\'on, pero no debe olvidarse que est\'a presente en el
hamiltoniano completo.

\section{\'Atomo de hidr\'ogeno: coordenadas de centro de masa}

Para un \'atomo de hidr\'ogeno se tienen dos cargas: electr\'on y prot\'on. En
la convenci\'on de las notas manuscritas, el t\'ermino lineal en el campo queda
\begin{equation}
  H_{\mathrm{int}}
  =
  -\frac{e}{m_e}
  \vb A(\vb r_e,t)\cdot\vb p_e
  +
  \frac{e}{m_p}
  \vb A(\vb r_p,t)\cdot\vb p_p .
\end{equation}
El signo relativo refleja que electr\'on y prot\'on tienen cargas opuestas.

Se introducen la coordenada de centro de masa y la coordenada relativa:
\begin{equation}
  \vb R
  =
  \frac{m_e\vb r_e+m_p\vb r_p}{M},
  \qquad
  \vb r
  =
  \vb r_e-\vb r_p,
  \qquad
  M=m_e+m_p.
\end{equation}
Para invertir estas relaciones, de \(\vb r=\vb r_e-\vb r_p\) se tiene
\(\vb r_p=\vb r_e-\vb r\). Sustituyendo en \(\vb R\),
\begin{align}
  \vb R
  &=
  \frac{m_e\vb r_e+m_p(\vb r_e-\vb r)}{M}                                      \\
  &=
  \vb r_e-\frac{m_p}{M}\vb r,
\end{align}
por lo que
\begin{equation}
  \vb r_e
  =
  \vb R+\frac{m_p}{M}\vb r.
\end{equation}
De manera an\'aloga,
\begin{equation}
  \vb r_p
  =
  \vb R-\frac{m_e}{M}\vb r.
\end{equation}

Los momentos conjugados convenientes son
\begin{equation}
  \vb P=\vb p_e+\vb p_p,
  \qquad
  \vb p
  =
  \frac{m_p}{M}\vb p_e
  -
  \frac{m_e}{M}\vb p_p,
\end{equation}
donde \(\vb P\) es el momento total y \(\vb p\) es el momento relativo. La masa
reducida es
\begin{equation}
  \mu
  =
  \frac{m_em_p}{M}.
\end{equation}
Para despejar \(\vb p_e\), se usa \(\vb p_p=\vb P-\vb p_e\):
\begin{align}
  \vb p
  &=
  \frac{m_p}{M}\vb p_e
  -
  \frac{m_e}{M}(\vb P-\vb p_e)                                                 \\
  &=
  \left(\frac{m_p+m_e}{M}\right)\vb p_e
  -
  \frac{m_e}{M}\vb P                                                           \\
  &=
  \vb p_e-\frac{m_e}{M}\vb P.
\end{align}
As\'i,
\begin{equation}
  \vb p_e
  =
  \frac{m_e}{M}\vb P+\vb p,
  \qquad
  \vb p_p
  =
  \frac{m_p}{M}\vb P-\vb p.
\end{equation}
La energ\'ia cin\'etica tambi\'en se separa. Sustituyendo,
\begin{align}
  \frac{\vb p_e^{\,2}}{2m_e}
  +
  \frac{\vb p_p^{\,2}}{2m_p}
  &=
  \frac{1}{2m_e}
  \left(
    \frac{m_e^2}{M^2}\vb P^{\,2}
    +
    \frac{2m_e}{M}\vb P\cdot\vb p
    +
    \vb p^{\,2}
  \right)                                                                     \\
  &\quad
  +
  \frac{1}{2m_p}
  \left(
    \frac{m_p^2}{M^2}\vb P^{\,2}
    -
    \frac{2m_p}{M}\vb P\cdot\vb p
    +
    \vb p^{\,2}
  \right).
\end{align}
Los t\'erminos cruzados se cancelan. Los t\'erminos restantes dan
\begin{equation}
  \frac{\vb p_e^{\,2}}{2m_e}
  +
  \frac{\vb p_p^{\,2}}{2m_p}
  =
  \frac{\vb P^{\,2}}{2M}
  +
  \frac{\vb p^{\,2}}{2\mu}.
\end{equation}
Esta separaci\'on es esencial: \(\vb P\) describe el movimiento del centro de
masa y \(\vb p\) describe la din\'amica interna del \'atomo.

\section{Aproximaci\'on dipolar en la forma \texorpdfstring{\(\vb A\cdot\vb p\)}{A.p}}

Las longitudes de onda \'opticas son mucho mayores que el tama\~no at\'omico.
T\'ipicamente,
\begin{equation}
  \lambda\sim 500\,\mathrm{nm},
  \qquad
  a_0\sim 0.05\,\mathrm{nm},
  \qquad
  a_0\ll\lambda .
\end{equation}
Por ello, el potencial vectorial casi no cambia a lo largo del \'atomo. Para
electr\'on y prot\'on,
\begin{equation}
  \vb A(\vb r_e,t)
  \simeq
  \vb A(\vb r_p,t)
  \simeq
  \vb A(\vb R,t).
\end{equation}
La estimaci\'on se obtiene expandiendo alrededor de \(\vb R\). Si
\(\vb r_e=\vb R+\delta\vb r\), entonces
\begin{equation}
  \vb A(\vb R+\delta\vb r,t)
  \simeq
  \vb A(\vb R,t)
  +
  (\delta\vb r\cdot\nabla_{\vb R})\vb A(\vb R,t)+\cdots .
\end{equation}
Para una onda \'optica,
\begin{equation}
  \vb A(\vb r,t)
  =
  \vb A_0\cos(\vb k\cdot\vb r-\omega t),
  \qquad
  |\vb k|=\frac{2\pi}{\lambda}.
\end{equation}
Cada derivada espacial introduce un factor del orden de \(|\vb k|\), de modo
que
\begin{equation}
  \frac{
  |(\delta\vb r\cdot\nabla_{\vb R})\vb A(\vb R,t)|
  }{
  |\vb A(\vb R,t)|
  }
  \sim
  |\delta\vb r|\,|\vb k|
  \lesssim
  |\vb r|\,\frac{2\pi}{\lambda}
  =
  2\pi\frac{\text{tama\~no at\'omico}}{\text{longitud de onda}}.
\end{equation}
Como esta raz\'on es mucho menor que uno, se conserva solamente el primer
t\'ermino.

Sustituyendo esta aproximaci\'on en el acoplamiento lineal,
\begin{align}
  H_{\mathrm{int}}
  &\simeq
  \vb A(\vb R,t)\cdot
  \left(
    -\frac{e}{m_e}\vb p_e
    +
    \frac{e}{m_p}\vb p_p
  \right)                                                                     \\
  &=
  \vb A(\vb R,t)\cdot
  \left[
    -\frac{e}{m_e}
    \left(
      \frac{m_e}{M}\vb P+\vb p
    \right)
    +
    \frac{e}{m_p}
    \left(
      \frac{m_p}{M}\vb P-\vb p
    \right)
  \right]                                                                     \\
  &=
  \vb A(\vb R,t)\cdot
  \left(
    -\frac{e}{M}\vb P
    -\frac{e}{m_e}\vb p
    +\frac{e}{M}\vb P
    -\frac{e}{m_p}\vb p
  \right).
\end{align}
Los t\'erminos proporcionales al momento total \(\vb P\) se cancelan. Queda
\begin{equation}
  H_{\vb A\cdot\vb p}
  =
  -e
  \left(
    \frac{1}{m_e}+\frac{1}{m_p}
  \right)
  \vb A(\vb R,t)\cdot\vb p.
\end{equation}
Como
\begin{equation}
  \frac{1}{\mu}
  =
  \frac{1}{m_e}+\frac{1}{m_p},
\end{equation}
se obtiene
\begin{equation}
  \boxed{
  H_{\vb A\cdot\vb p}
  =
  -\frac{e}{\mu}
  \vb A(\vb R,t)\cdot\vb p .
  }
\end{equation}
Esta es la forma potencial vectorial--momento de la interacci\'on en la
aproximaci\'on dipolar.

\section{Hamiltoniano hidrogenoide en el campo}

El potencial interno del hidr\'ogeno es el potencial coulombiano
\begin{equation}
  V(\vb r)
  =
  -\frac{1}{4\pi\varepsilon_0}
  \frac{e^2}{|\vb r|}.
\end{equation}
El hamiltoniano, antes de hacer la aproximaci\'on dipolar, puede organizarse
como la suma de las energ\'ias cin\'eticas de electr\'on y prot\'on, el
potencial coulombiano y los acoplamientos con el campo:
\begin{equation}
  H
  =
  H_e+H_p+V(\vb r_e-\vb r_p),
\end{equation}
con
\begin{align}
  H_e
  &=
  \frac{\vb p_e^{\,2}}{2m_e}
  -
  \frac{e}{m_e}
  \vb A(\vb r_e,t)\cdot\vb p_e
  +
  \frac{e^2}{2m_e}
  \vb A^2(\vb r_e,t),                                                         \\
  H_p
  &=
  \frac{\vb p_p^{\,2}}{2m_p}
  +
  \frac{e}{m_p}
  \vb A(\vb r_p,t)\cdot\vb p_p
  +
  \frac{e^2}{2m_p}
  \vb A^2(\vb r_p,t).
\end{align}
Al pasar a \(\vb R,\vb r,\vb P,\vb p\), la parte cin\'etica libre se convierte
en
\begin{equation}
  \frac{\vb P^{\,2}}{2M}+\frac{\vb p^{\,2}}{2\mu}.
\end{equation}
La parte de interacci\'on contiene dos tipos de t\'erminos lineales: unos
acoplan el campo al momento relativo \(\vb p\), y otros acoplan el campo al
momento de centro de masa \(\vb P\). En orden dipolar, como
\(\vb A(\vb r_e,t)\simeq\vb A(\vb r_p,t)\simeq\vb A(\vb R,t)\), los t\'erminos
con \(\vb P\) se cancelan y queda el acoplamiento relativo derivado arriba.

Conservando tambi\'en el t\'ermino cuadr\'atico en el campo, el hamiltoniano en
orden dipolar es
\begin{equation}
  H^{(0)}
  =
  \frac{\vb P^{\,2}}{2M}
  +
  \frac{\vb p^{\,2}}{2\mu}
  -
  \frac{1}{4\pi\varepsilon_0}
  \frac{e^2}{|\vb r|}
  -
  \frac{e}{\mu}\vb A(\vb R,t)\cdot\vb p
  +
  \frac{e^2}{2\mu}\vb A^2(\vb R,t).
\end{equation}
Equivalentemente,
\begin{equation}
  H^{(0)}
  =
  \frac{\vb P^{\,2}}{2M}
  +
  \frac{1}{2\mu}
  \left[
    \vb p-e\vb A(\vb R,t)
  \right]^2
  -
  \frac{1}{4\pi\varepsilon_0}
  \frac{e^2}{|\vb r|}.
\end{equation}
Si el centro de masa se trata aparte, se identifica
\begin{equation}
  H_{\mathrm{atom}}
  =
  \frac{\vb p^{\,2}}{2\mu}
  -
  \frac{1}{4\pi\varepsilon_0}
  \frac{e^2}{|\vb r|},
\end{equation}
y
\begin{equation}
  H_{\mathrm{int}}
  =
  -\frac{e}{\mu}\vb A(\vb R,t)\cdot\vb p.
\end{equation}
El t\'ermino
\begin{equation}
  U(\vb R)
  =
  \frac{e^2}{2\mu}\vb A^2(\vb R,t)
\end{equation}
act\'ua como potencial ponderomotriz para el centro de masa. Es importante en
contextos de campos intensos, plasmas, \'atomos en campos fuertes y l\'aseres
de electrones libres.

\section{Invariancia local de gauge y acoplamiento m\'inimo}

La forma de acoplamiento m\'inimo tambi\'en puede entenderse como consecuencia
de la invariancia local de gauge. Consid\'erese primero la ecuaci\'on de
Schr\"odinger sin campo electromagn\'etico,
\begin{equation}
  \ii\hbar\frac{\partial\psi}{\partial t}
  =
  \left[
    \frac{1}{2m}
    \left(
      \frac{\hbar}{\ii}\nabla
    \right)^2
    +
    U(\vb r)
  \right]\psi.
\end{equation}
La cantidad observable directa asociada a la funci\'on de onda es
\(|\psi|^2\). Por tanto, una fase global
\begin{equation}
  \tilde\psi(\vb r,t)
  =
  \ee^{\ii\alpha}\psi(\vb r,t),
\end{equation}
con \(\alpha\) constante, no cambia la f\'isica. Adem\'as, como la fase no
depende de \(\vb r\) ni de \(t\), \(\tilde\psi\) satisface la misma ecuaci\'on
de Schr\"odinger.

La pregunta de Weyl fue m\'as profunda: \textquestiondown qu\'e ocurre si la
fase puede elegirse localmente, punto por punto? Se toma
\begin{equation}
  \tilde\psi(\vb r,t)
  =
  \ee^{\frac{\ii}{\hbar}e\Lambda(\vb r,t)}
  \psi(\vb r,t).
\end{equation}
Al derivar respecto al tiempo,
\begin{align}
  \ii\hbar\frac{\partial\tilde\psi}{\partial t}
  &=
  \ee^{\frac{\ii}{\hbar}e\Lambda}
  \left(
    \ii\hbar\frac{\partial\psi}{\partial t}
    -
    e\frac{\partial\Lambda}{\partial t}\psi
  \right).
\end{align}
Al derivar espacialmente,
\begin{align}
  \frac{\hbar}{\ii}\nabla\tilde\psi
  &=
  \ee^{\frac{\ii}{\hbar}e\Lambda}
  \left[
    \frac{\hbar}{\ii}\nabla\psi
    +
    e(\nabla\Lambda)\psi
  \right].
\end{align}
Por tanto,
\begin{equation}
  \left(
    \frac{\hbar}{\ii}\nabla
  \right)^2\tilde\psi
  =
  \ee^{\frac{\ii}{\hbar}e\Lambda}
  \left(
    \frac{\hbar}{\ii}\nabla+e\nabla\Lambda
  \right)^2\psi.
\end{equation}
La ecuaci\'on libre no mantiene su forma si la fase depende de
\((\vb r,t)\). Para restaurar la invariancia se introducen potenciales que se
transforman como
\begin{equation}
  \vb A'
  =
  \vb A+\nabla\Lambda,
  \qquad
  \phi'
  =
  \phi-\frac{\partial\Lambda}{\partial t}.
\end{equation}
Entonces
\begin{equation}
  \vb B'
  =
  \nabla\times\vb A'
  =
  \nabla\times\vb A
  +
  \nabla\times\nabla\Lambda
  =
  \vb B,
\end{equation}
y
\begin{align}
  \vb E'
  &=
  -\nabla\phi'
  -
  \frac{\partial\vb A'}{\partial t}                                           \\
  &=
  -\nabla\phi
  +
  \nabla\frac{\partial\Lambda}{\partial t}
  -
  \frac{\partial\vb A}{\partial t}
  -
  \frac{\partial}{\partial t}\nabla\Lambda                                    \\
  &=
  \vb E.
\end{align}
Los campos f\'isicos no cambian. La ecuaci\'on que s\'i es covariante bajo
transformaciones locales de fase es
\begin{equation}
  \ii\hbar\frac{\partial\psi}{\partial t}
  =
  \left[
    \frac{1}{2m}
    \left(
      \frac{\hbar}{\ii}\nabla-e\vb A(\vb r,t)
    \right)^2
    +
    U(\vb r)
    +
    e\phi(\vb r,t)
  \right]\psi .
\end{equation}
As\'i, la invariancia local de gauge obliga a reemplazar el momento can\'onico
por el momento mec\'anico \(\vb p-e\vb A\). Esta es la raz\'on estructural del
acoplamiento m\'inimo.

\section{Forma campo el\'ectrico--dipolo}

En muchas aplicaciones de \'optica cu\'antica conviene escribir la interacci\'on
en t\'erminos del campo el\'ectrico y del dipolo del \'atomo. El momento dipolar
asociado a la coordenada relativa es
\begin{equation}
  \vb d=e\vb r.
\end{equation}
Como el campo el\'ectrico tambi\'en var\'ia muy poco a lo largo del \'atomo,
la energ\'ia de un dipolo en un campo uniforme es
\begin{equation}
  \boxed{
  H_{\vb r\cdot\vb E}
  =
  -e\vb r\cdot\vb E(\vb R,t)
  =
  -\vb d\cdot\vb E(\vb R,t).
  }
\end{equation}
En esta forma, el campo acopla directamente los dos centros de carga mediante
la separaci\'on \(\vb r\). Es la forma que se usar\'a para construir el modelo
de dos niveles porque sus elementos de matriz se interpretan naturalmente como
dipolos de transici\'on.

Si el centro de masa est\'a fijo, el hamiltoniano dipolar queda
\begin{equation}
  \tilde H^{(0)}
  =
  \frac{\vb P^{\,2}}{2M}
  +
  \frac{\vb p^{\,2}}{2\mu}
  +
  V(\vb r)
  -
  e\vb r\cdot\vb E(\vb R,t).
\end{equation}
Cuando el movimiento del centro de masa es relevante, la equivalencia entre las
formas \(\vb A\cdot\vb p\) y \(\vb r\cdot\vb E\) requiere cuidado adicional.
En particular, aparecen correcciones tipo R\"ontgen, asociadas al hecho de que
un observador en movimiento en presencia de un campo magn\'etico percibe un
campo el\'ectrico efectivo. En esta nota se conserva la aproximaci\'on usual
con centro de masa fijo o tratado de forma separada.

\section{Equivalencia cl\'asica de las dos formas}

La equivalencia cl\'asica puede verse usando el hecho de que dos lagrangianos
que difieren por una derivada total generan las mismas ecuaciones de movimiento.
La acci\'on es
\begin{equation}
  S
  =
  \int_{t_1}^{t_2}
  L(\vb r,\dot{\vb r},t)\,\dd t.
\end{equation}
Si se define
\begin{equation}
  \tilde L
  =
  L+
  \frac{\dd}{\dd t}F(\vb r,t),
\end{equation}
entonces
\begin{equation}
  \tilde S
  =
  S+
  F(\vb r(t_2),t_2)
  -
  F(\vb r(t_1),t_1).
\end{equation}
Como en el principio variacional los extremos est\'an fijos, la variaci\'on de
los t\'erminos de borde es cero. Por tanto, \(L\) y \(\tilde L\) describen la
misma din\'amica.

Con el centro de masa fijo, el hamiltoniano de acoplamiento m\'inimo es
\begin{equation}
  H^{(0)}
  =
  \frac{1}{2\mu}
  \left[
    \vb p-e\vb A(\vb R,t)
  \right]^2
  +
  V(\vb r).
\end{equation}
La relaci\'on entre velocidad y momento se obtiene de
\begin{equation}
  \dot{\vb r}
  =
  \frac{\partial H^{(0)}}{\partial\vb p}
  =
  \frac{1}{\mu}
  \left[
    \vb p-e\vb A(\vb R,t)
  \right],
\end{equation}
por lo que
\begin{equation}
  \vb p
  =
  \mu\dot{\vb r}
  +
  e\vb A(\vb R,t).
\end{equation}
La transformada inversa de Legendre da
\begin{align}
  L^{(0)}
  &=
  \vb p\cdot\dot{\vb r}-H^{(0)}                                               \\
  &=
  \frac{\mu}{2}\dot{\vb r}^{\,2}
  +
  e\dot{\vb r}\cdot\vb A(\vb R,t)
  -
  V(\vb r).
\end{align}
Ahora se resta la derivada total
\begin{equation}
  e\frac{\dd}{\dd t}\left[\vb r\cdot\vb A(\vb R,t)\right].
\end{equation}
Como \(\vb R\) se considera fijo,
\begin{equation}
  \frac{\dd}{\dd t}\left[\vb r\cdot\vb A(\vb R,t)\right]
  =
  \dot{\vb r}\cdot\vb A(\vb R,t)
  +
  \vb r\cdot
  \frac{\partial\vb A(\vb R,t)}{\partial t}.
\end{equation}
Entonces
\begin{align}
  \tilde L^{(0)}
  &=
  L^{(0)}
  -
  e\frac{\dd}{\dd t}\left[\vb r\cdot\vb A(\vb R,t)\right]                     \\
  &=
  \frac{\mu}{2}\dot{\vb r}^{\,2}
  +
  e\dot{\vb r}\cdot\vb A
  -
  V(\vb r)
  -
  e\dot{\vb r}\cdot\vb A
  -
  e\vb r\cdot\frac{\partial\vb A}{\partial t}                                \\
  &=
  \frac{\mu}{2}\dot{\vb r}^{\,2}
  -
  V(\vb r)
  -
  e\vb r\cdot\frac{\partial\vb A}{\partial t}.
\end{align}
En el gauge usado para el campo de radiaci\'on,
\begin{equation}
  \vb E(\vb R,t)
  =
  -\frac{\partial\vb A(\vb R,t)}{\partial t},
\end{equation}
por lo que
\begin{equation}
  \tilde L^{(0)}
  =
  \frac{\mu}{2}\dot{\vb r}^{\,2}
  -
  V(\vb r)
  +
  e\vb r\cdot\vb E(\vb R,t).
\end{equation}
El hamiltoniano correspondiente es
\begin{equation}
  \tilde H^{(0)}
  =
  \frac{\vb p^{\,2}}{2\mu}
  +
  V(\vb r)
  -
  e\vb r\cdot\vb E(\vb R,t).
\end{equation}
As\'i se obtiene la forma el\'ectrico--dipolo a partir de la forma
\(\vb A\cdot\vb p\).

\section{Equivalencia cu\'antica por transformaci\'on de gauge}

La equivalencia cu\'antica se expresa mediante una transformaci\'on de fase.
Consid\'erese la ecuaci\'on de Schr\"odinger en la forma de acoplamiento
m\'inimo,
\begin{equation}
  \ii\hbar\frac{\partial\tilde\psi(\vb r,t)}{\partial t}
  =
  \left[
    \frac{1}{2\mu}
    \left(
      \vb p-e\vb A(\vb R,t)
    \right)^2
    +
    V(\vb r)
  \right]\tilde\psi(\vb r,t).
\end{equation}
Se introduce la fase
\begin{equation}
  \tilde\psi(\vb r,t)
  =
  \ee^{\frac{\ii}{\hbar}e\vb r\cdot\vb A(\vb R,t)}
  \psi(\vb r,t).
\end{equation}
Esta transformaci\'on corresponde a elegir
\begin{equation}
  \Lambda(\vb r,t)
  =
  \vb r\cdot\vb A(\vb R,t).
\end{equation}
Como la derivada se toma respecto a la coordenada interna \(\vb r\),
\begin{equation}
  \nabla_{\vb r}\Lambda
  =
  \vb A(\vb R,t).
\end{equation}
Con la elecci\'on de signo adecuada para pasar de \(\tilde\psi\) a \(\psi\),
el potencial vectorial transformado es
\begin{equation}
  \vb A'
  =
  \vb A-\nabla_{\vb r}\Lambda
  =
  \vb A-\vb A
  =
  0.
\end{equation}
El potencial escalar transformado es
\begin{equation}
  \phi'
  =
  \phi+\frac{\partial\Lambda}{\partial t}.
\end{equation}
Tomando \(\phi=0\) para el campo de radiaci\'on,
\begin{equation}
  \phi'
  =
  \vb r\cdot
  \frac{\partial\vb A(\vb R,t)}{\partial t}
  =
  -\vb r\cdot\vb E(\vb R,t).
\end{equation}
Por tanto, \(\psi\) satisface
\begin{equation}
  \ii\hbar\frac{\partial\psi(\vb r,t)}{\partial t}
  =
  \left[
    \frac{\vb p^{\,2}}{2\mu}
    +
    V(\vb r)
    -
    e\vb r\cdot\vb E(\vb R,t)
  \right]\psi(\vb r,t).
\end{equation}
La conclusi\'on es que si \(\tilde\psi\) obedece la ecuaci\'on con
\(\vb A\cdot\vb p\), entonces \(\psi\), que es la funci\'on de onda
transformada por gauge, obedece la ecuaci\'on con \(\vb r\cdot\vb E\).

\section{Elementos de matriz de las dos interacciones}

Aunque los dos hamiltonianos son equivalentes, sus elementos de matriz entre
los mismos eigenestados del \'atomo libre no tienen por qu\'e coincidir. Esta
observaci\'on es importante porque evita una confusi\'on com\'un: la equivalencia
de gauge se refiere a la din\'amica completa, incluyendo la transformaci\'on de
los estados, no a comparar directamente cada operador en la misma base sin
transformar.

Los hamiltonianos de interacci\'on son
\begin{equation}
  \hat H_{\vb A\cdot\vb p}
  =
  -\frac{e}{\mu}
  \vb A(\vb R,t)\cdot\hat{\vb p},
  \qquad
  \hat H_{\vb r\cdot\vb E}
  =
  -e\hat{\vb r}\cdot\vb E(\vb R,t).
\end{equation}
Aqu\'i \(\vb A\) y \(\vb E\) se tratan como campos cl\'asicos para discutir los
elementos de matriz at\'omicos. Si \(\ket j\) y \(\ket{j'}\) son eigenestados
del hamiltoniano at\'omico libre,
\begin{equation}
  \hat H_{\mathrm{atom}}\ket j
  =
  E_j\ket j,
\end{equation}
entonces
\begin{align}
  \bra j\hat H_{\vb A\cdot\vb p}\ket{j'}
  &=
  -\frac{e}{\mu}
  \vb A(\vb R,t)\cdot
  \bra j\hat{\vb p}\ket{j'},                                                   \\
  \bra j\hat H_{\vb r\cdot\vb E}\ket{j'}
  &=
  -e
  \vb E(\vb R,t)\cdot
  \bra j\hat{\vb r}\ket{j'}.
\end{align}
Para relacionar \(\hat{\vb p}\) con \(\hat{\vb r}\), se usa la ecuaci\'on de
Heisenberg con el hamiltoniano libre:
\begin{equation}
  \frac{\hat{\vb p}}{\mu}
  =
  \frac{\dd\hat{\vb r}}{\dd t}
  =
  \frac{\ii}{\hbar}
  \left[
    \hat H_{\mathrm{atom}},
    \hat{\vb r}
  \right].
\end{equation}
Tomando elementos de matriz,
\begin{align}
  \frac{1}{\mu}
  \bra j\hat{\vb p}\ket{j'}
  &=
  \frac{\ii}{\hbar}
  \bra j
  \hat H_{\mathrm{atom}}\hat{\vb r}
  -
  \hat{\vb r}\hat H_{\mathrm{atom}}
  \ket{j'}                                                                    \\
  &=
  \frac{\ii}{\hbar}
  \left(
    E_j-E_{j'}
  \right)
  \bra j\hat{\vb r}\ket{j'}.
\end{align}
Definiendo
\begin{equation}
  \omega_{jj'}
  =
  \frac{E_j-E_{j'}}{\hbar},
\end{equation}
se obtiene
\begin{equation}
  \frac{1}{\mu}
  \bra j\hat{\vb p}\ket{j'}
  =
  \ii\omega_{jj'}
  \bra j\hat{\vb r}\ket{j'}.
\end{equation}
Entonces
\begin{equation}
  \bra j\hat H_{\vb A\cdot\vb p}\ket{j'}
  =
  -\ii e\omega_{jj'}
  \vb A(\vb R,t)\cdot
  \bra j\hat{\vb r}\ket{j'}.
\end{equation}
Esto no es, en general, igual a
\begin{equation}
  -e\vb E(\vb R,t)\cdot\bra j\hat{\vb r}\ket{j'}.
\end{equation}
Por tanto,
\begin{equation}
  \bra j\hat H_{\vb A\cdot\vb p}\ket{j'}
  \neq
  \bra j\hat H_{\vb r\cdot\vb E}\ket{j'}.
\end{equation}
La diferencia no contradice la equivalencia de las dos descripciones. Solo
recuerda que al cambiar de gauge tambi\'en cambia la representaci\'on de los
estados.

\section{Subsistemas, interacci\'on y entrelazamiento}

En la teor\'ia cu\'antica completa, el estado de \'atomo m\'as campo satisface
\begin{equation}
  \ii\hbar\frac{\dd}{\dd t}\ket{\Psi}
  =
  \hat H\ket{\Psi}.
\end{equation}
Antes de despreciar el movimiento del centro de masa, el estado total puede
pensarse como un producto tensorial de tres sectores:
\begin{equation}
  \ket{\Psi}
  =
  \ket{\psi_{\mathrm{cm}}}
  \ket{\psi_{\mathrm{atom}}}
  \ket{\psi_{\mathrm{campo}}}.
\end{equation}
Los operadores \(\hat{\vb R},\hat{\vb P}\) act\'uan sobre el centro de masa; los
operadores \(\hat{\vb r},\hat{\vb p}\) act\'uan sobre los grados internos del
\'atomo; y los operadores \(\hat{\vb A},\hat{\vb E},\hat{\vb B}\) contienen
\(\hat a,\adag\) y las funciones de modo \(\vb u_l(\vb R)\), por lo que act\'uan
sobre el campo y pueden depender de la posici\'on del \'atomo.

Aunque el estado inicial sea separable, una interacci\'on como
\(\hat H_{\vb r\cdot\vb E}\) multiplica operadores at\'omicos por operadores del
campo. Por eso la evoluci\'on genera, en general, un estado entrelazado. Esta es
la base f\'isica del modelo de Jaynes--Cummings--Paul: el \'atomo no solo
intercambia energ\'ia con el modo de la cavidad, sino que se correlaciona
cu\'anticamente con el n\'umero y la fase del campo.

\section{Modelo simple de interacci\'on \'atomo--campo}

Para llegar a un modelo manejable se hacen tres simplificaciones:
\begin{enumerate}
  \item se conservan solo dos niveles at\'omicos, uno excitado \(\ket a\) y uno
  base \(\ket b\);
  \item se conserva un solo modo del campo de la cavidad;
  \item inicialmente se permite que el acoplamiento dependa de la posici\'on del
  centro de masa, aunque al ignorar este movimiento se obtiene el modelo usual
  de Jaynes--Cummings--Paul.
\end{enumerate}

Las energ\'ias de los dos estados at\'omicos son
\begin{equation}
  E_a=\hbar\omega_a,
  \qquad
  E_b=\hbar\omega_b,
\end{equation}
y la frecuencia de transici\'on es
\begin{equation}
  \omega=\omega_a-\omega_b
  =
  \frac{E_a-E_b}{\hbar}.
\end{equation}
La identidad en el subespacio de dos niveles es
\begin{equation}
  \mathbb{1}
  =
  \ket a\bra a+\ket b\bra b.
\end{equation}
Para cualquier operador at\'omico \(\hat O\),
\begin{equation}
  \hat O
  =
  \mathbb{1}\hat O\mathbb{1}
  =
  \sum_{j=a,b}\sum_{j'=a,b}
  \ket j
  \bra j\hat O\ket{j'}
  \bra{j'}.
\end{equation}
Aplicando esto al hamiltoniano at\'omico,
\begin{align}
  \hat H_{\mathrm{atom}}
  &=
  \mathbb{1}\hat H_{\mathrm{atom}}\mathbb{1}                                  \\
  &=
  E_a\ket a\bra a
  +
  E_b\ket b\bra b,
\end{align}
porque \(\ket a\) y \(\ket b\) son eigenestados de energ\'ia, de modo que los
elementos fuera de la diagonal se anulan.

En la base
\begin{equation}
  \ket a=
  \begin{pmatrix}
    1\\0
  \end{pmatrix},
  \qquad
  \ket b=
  \begin{pmatrix}
    0\\1
  \end{pmatrix},
\end{equation}
el hamiltoniano at\'omico es
\begin{equation}
  \hat H_{\mathrm{atom}}
  =
  \begin{pmatrix}
    E_a & 0\\
    0 & E_b
  \end{pmatrix}.
\end{equation}
Se separa una constante de energ\'ia:
\begin{align}
  \hat H_{\mathrm{atom}}
  &=
  \begin{pmatrix}
    \frac{E_a+E_b}{2}+\frac{E_a-E_b}{2} & 0\\
    0 & \frac{E_a+E_b}{2}-\frac{E_a-E_b}{2}
  \end{pmatrix}                                                               \\
  &=
  E\mathbb{1}
  +
  \frac{\hbar\omega}{2}
  \begin{pmatrix}
    1&0\\
    0&-1
  \end{pmatrix},
\end{align}
donde
\begin{equation}
  E=\frac{E_a+E_b}{2}.
\end{equation}
El t\'ermino \(E\mathbb{1}\) solo produce una fase global, por lo que se omite.
Definiendo
\begin{equation}
  \hat\sigma_z
  =
  \ket a\bra a-\ket b\bra b
  =
  \begin{pmatrix}
    1&0\\
    0&-1
  \end{pmatrix},
\end{equation}
queda
\begin{equation}
  \boxed{
  \hat H_{\mathrm{atom}}
  =
  \frac{\hbar\omega}{2}\hat\sigma_z .
  }
\end{equation}

\section{Operador dipolar en el subespacio de dos niveles}

El operador de posici\'on \(\hat{\vb r}\) es impar bajo paridad:
\(\vb r\to-\vb r\). Por ello, si los estados at\'omicos tienen paridad definida,
los elementos diagonales se anulan. En detalle,
\begin{equation}
  \bra j\hat{\vb r}\ket j
  =
  \int \dd^3r\,
  |\psi_j(\vb r)|^2\vb r.
\end{equation}
Haciendo el cambio de variable \(\vb r=-\vb u\), se tiene
\(\dd^3r=\dd^3u\) y
\begin{align}
  \bra j\hat{\vb r}\ket j
  &=
  -\int \dd^3u\,
  |\psi_j(-\vb u)|^2\vb u.
\end{align}
Si \(\ket j\) tiene paridad definida, entonces
\(|\psi_j(-\vb u)|^2=|\psi_j(\vb u)|^2\). Por tanto, la integral es igual a su
negativo:
\begin{equation}
  \bra j\hat{\vb r}\ket j
  =
  -\bra j\hat{\vb r}\ket j,
\end{equation}
y necesariamente
\begin{equation}
  \bra j\hat{\vb r}\ket j=0.
\end{equation}

Los elementos no diagonales s\'i pueden ser distintos de cero. Se define el
dipolo de transici\'on
\begin{equation}
  \vb d
  =
  e\bra a\hat{\vb r}\ket b
  =
  e\int \dd^3r\,\psi_a^*(\vb r)\,\vb r\,\psi_b(\vb r),
\end{equation}
y su conjugado
\begin{equation}
  \vb d^{\,*}
  =
  e\bra b\hat{\vb r}\ket a.
\end{equation}
Entonces
\begin{align}
  e\hat{\vb r}
  &=
  e\mathbb{1}\hat{\vb r}\mathbb{1}                                            \\
  &=
  e\ket a\bra a\hat{\vb r}\ket a\bra a
  +
  e\ket a\bra a\hat{\vb r}\ket b\bra b                                      \\
  &\quad
  +
  e\ket b\bra b\hat{\vb r}\ket a\bra a
  +
  e\ket b\bra b\hat{\vb r}\ket b\bra b                                      \\
  &=
  \vb d\,\ket a\bra b
  +
  \vb d^{\,*}\ket b\bra a.
\end{align}
Este operador no desplaza al \'atomo dentro del mismo nivel. Lo que hace es
conectar los dos niveles:
\begin{equation}
  e\hat{\vb r}\ket b
  =
  \vb d\ket a,
  \qquad
  e\hat{\vb r}\ket a
  =
  \vb d^{\,*}\ket b.
\end{equation}
Se introducen los operadores de subida y bajada
\begin{equation}
  \splus
  =
  \ket a\bra b
  =
  \begin{pmatrix}
    0&1\\
    0&0
  \end{pmatrix},
  \qquad
  \sminus
  =
  \ket b\bra a
  =
  \begin{pmatrix}
    0&0\\
    1&0
  \end{pmatrix}.
\end{equation}
El operador \(\splus\) crea una excitaci\'on at\'omica porque lleva
\(\ket b\) a \(\ket a\), mientras que \(\sminus\) destruye una excitaci\'on
at\'omica porque lleva \(\ket a\) a \(\ket b\). En esta notaci\'on,
\begin{equation}
  \boxed{
  e\hat{\vb r}
  =
  \vb d\,\splus+\vb d^{\,*}\sminus .
  }
\end{equation}

\section{Campo de un solo modo y frecuencia de Rabi del vac\'io}

Para un solo modo de la cavidad, de frecuencia \(\Omega\), el hamiltoniano del
campo se escribe, omitiendo la energ\'ia de punto cero,
\begin{equation}
  \hat H_{\mathrm{campo}}
  =
  \hbar\Omega\,\adag\hat a.
\end{equation}
El campo el\'ectrico de ese modo, evaluado en la posici\'on del \'atomo, se toma
como
\begin{equation}
  \hat{\vb E}(\vb R,t)
  =
  \Ezero\,\vb u(\vb R)\,\ii(\hat a-\adag),
\end{equation}
donde
\begin{equation}
  \Ezero
  =
  \sqrt{\frac{\hbar\Omega}{2\varepsilon_0V}}
\end{equation}
es la amplitud del campo el\'ectrico de un fot\'on, \(V\) es el volumen efectivo
del modo y \(\vb u(\vb R)\) contiene la forma espacial y polarizaci\'on del
modo.

La interacci\'on el\'ectrico--dipolo es
\begin{align}
  \hat H_{\vb r\cdot\vb E}
  &=
  -e\hat{\vb r}\cdot\hat{\vb E}(\vb R,t)                                      \\
  &=
  -
  \left(
    \vb d\,\splus+\vb d^{\,*}\sminus
  \right)
  \cdot
  \Ezero\vb u(\vb R)\,\ii(\hat a-\adag)                                      \\
  &=
  -\ii\Ezero
  \left[
    \vb d\cdot\vb u(\vb R)\,\splus
    +
    \vb d^{\,*}\cdot\vb u(\vb R)\,\sminus
  \right]
  (\hat a-\adag).
\end{align}
Se define la frecuencia de Rabi del vac\'io dependiente de la posici\'on:
\begin{equation}
  g(\vb R)
  =
  \frac{\Ezero}{\hbar}
  \left|
    \vb d\cdot\vb u(\vb R)
  \right|.
\end{equation}
Si
\begin{equation}
  \vb d\cdot\vb u(\vb R)
  =
  \left|
    \vb d\cdot\vb u(\vb R)
  \right|
  \ee^{\ii\varphi},
\end{equation}
entonces
\begin{equation}
  \hat H_{\vb r\cdot\vb E}
  =
  \hbar g(\vb R)(-\ii)
  \left(
    \splus\ee^{\ii\varphi}
    +
    \sminus\ee^{-\ii\varphi}
  \right)
  (\hat a-\adag).
\end{equation}
Escogiendo la fase \(\varphi=\pi/2\), se tiene
\(\ee^{\ii\varphi}=\ii\) y \(\ee^{-\ii\varphi}=-\ii\). Por tanto,
\begin{equation}
  \hat H_{\vb r\cdot\vb E}
  =
  \hbar g(\vb R)
  (\splus-\sminus)
  (\hat a-\adag).
\end{equation}
Esta elecci\'on de fase no cambia la f\'isica; solo fija una convenci\'on para
que el acoplamiento quede real.

\section{Hamiltoniano total antes de la aproximaci\'on de onda rotante}

Reuniendo el movimiento de centro de masa, el campo, el \'atomo de dos niveles y
la interacci\'on, se obtiene
\begin{equation}
  \boxed{
  \hat H
  =
  \frac{\hat{\vb P}^{\,2}}{2M}
  +
  \hbar\Omega\,\adag\hat a
  +
  \frac{\hbar\omega}{2}\hat\sigma_z
  +
  \hbar g(\hat{\vb R})
  (\splus-\sminus)(\hat a-\adag).
  }
\end{equation}
Si por el momento se omite la dependencia espacial de \(g\), la interacci\'on se
desarrolla como
\begin{align}
  \hat H_{\vb r\cdot\vb E}
  &=
  \hbar g(\splus-\sminus)(\hat a-\adag)                                      \\
  &=
  \hbar g
  \left(
    \splus\hat a
    -
    \splus\adag
    -
    \sminus\hat a
    +
    \sminus\adag
  \right).
\end{align}
Los cuatro t\'erminos tienen interpretaciones distintas:
\begin{itemize}
  \item \(\splus\hat a\): el campo pierde un fot\'on y el \'atomo se excita;
  \item \(\sminus\adag\): el \'atomo decae y el campo gana un fot\'on;
  \item \(\splus\adag\): se crea una excitaci\'on at\'omica y tambi\'en un fot\'on;
  \item \(\sminus\hat a\): se destruye una excitaci\'on at\'omica y tambi\'en un
  fot\'on.
\end{itemize}
Los dos primeros conservan el n\'umero total de excitaciones. Los dos \'ultimos
son contra-rotantes: cambian al \'atomo y al campo en la misma direcci\'on
energ\'etica, por lo que oscilan r\'apidamente cuando el sistema est\'a cerca
de resonancia.

\section{Cuadro de interacci\'on}

Para justificar la aproximaci\'on de onda rotante de manera expl\'icita se pasa
al cuadro de interacci\'on con respecto a
\begin{equation}
  \hat H_0
  =
  \hbar\Omega\,\adag\hat a
  +
  \frac{\hbar\omega}{2}\hat\sigma_z.
\end{equation}
Se define
\begin{equation}
  \ket{\Psi_I(t)}
  =
  \ee^{\frac{\ii}{\hbar}\hat H_0t}
  \ket{\Psi(t)},
  \qquad
  \ket{\Psi(t)}
  =
  \ee^{-\frac{\ii}{\hbar}\hat H_0t}
  \ket{\Psi_I(t)}.
\end{equation}
Sustituyendo en la ecuaci\'on de Schr\"odinger,
\begin{equation}
  \ii\hbar\frac{\dd}{\dd t}\ket{\Psi(t)}
  =
  \left(
    \frac{\hat{\vb P}^{\,2}}{2M}
    +
    \hat H_0
    +
    \hat H_{\vb r\cdot\vb E}
  \right)
  \ket{\Psi(t)},
\end{equation}
se obtiene
\begin{equation}
  \ii\hbar\frac{\dd}{\dd t}\ket{\Psi_I(t)}
  =
  \left[
    \frac{\hat{\vb P}^{\,2}}{2M}
    +
    \hat H_{\vb r\cdot\vb E}^{(I)}(t)
  \right]
  \ket{\Psi_I(t)}.
\end{equation}
El t\'ermino de centro de masa permanece igual porque conmuta con
\(\hat H_0\). La interacci\'on en el cuadro de interacci\'on es
\begin{equation}
  \hat H_{\vb r\cdot\vb E}^{(I)}(t)
  =
  \ee^{\frac{\ii}{\hbar}\hat H_0t}
  \hat H_{\vb r\cdot\vb E}
  \ee^{-\frac{\ii}{\hbar}\hat H_0t}.
\end{equation}
Como los operadores del campo conmutan con los operadores at\'omicos, la
transformaci\'on se separa:
\begin{align}
  \hat H_{\vb r\cdot\vb E}^{(I)}(t)
  &=
  \hbar g\,
  \ee^{\frac{\ii}{2}\omega t\hat\sigma_z}
  (\splus-\sminus)
  \ee^{-\frac{\ii}{2}\omega t\hat\sigma_z}                                   \\
  &\quad\times
  \ee^{\ii\Omega t\adag\hat a}
  (\hat a-\adag)
  \ee^{-\ii\Omega t\adag\hat a}.
\end{align}

Para la parte at\'omica,
\begin{equation}
  \ee^{\frac{\ii}{2}\omega t\hat\sigma_z}
  =
  \begin{pmatrix}
    \ee^{\ii\omega t/2}&0\\
    0&\ee^{-\ii\omega t/2}
  \end{pmatrix}.
\end{equation}
Al multiplicar matrices,
\begin{align}
  \ee^{\frac{\ii}{2}\omega t\hat\sigma_z}
  \splus
  \ee^{-\frac{\ii}{2}\omega t\hat\sigma_z}
  &=
  \ee^{\ii\omega t}\splus,                                                     \\
  \ee^{\frac{\ii}{2}\omega t\hat\sigma_z}
  \sminus
  \ee^{-\frac{\ii}{2}\omega t\hat\sigma_z}
  &=
  \ee^{-\ii\omega t}\sminus.
\end{align}
Por tanto,
\begin{equation}
  \ee^{\frac{\ii}{2}\omega t\hat\sigma_z}
  (\splus-\sminus)
  \ee^{-\frac{\ii}{2}\omega t\hat\sigma_z}
  =
  \ee^{\ii\omega t}\splus
  -
  \ee^{-\ii\omega t}\sminus.
\end{equation}

Para el campo,
\begin{equation}
  \ee^{\ii\Omega t\adag\hat a}
  \hat a
  \ee^{-\ii\Omega t\adag\hat a}
  =
  \ee^{-\ii\Omega t}\hat a,
\end{equation}
y
\begin{equation}
  \ee^{\ii\Omega t\adag\hat a}
  \adag
  \ee^{-\ii\Omega t\adag\hat a}
  =
  \ee^{\ii\Omega t}\adag.
\end{equation}
Entonces
\begin{equation}
  \ee^{\ii\Omega t\adag\hat a}
  (\hat a-\adag)
  \ee^{-\ii\Omega t\adag\hat a}
  =
  \ee^{-\ii\Omega t}\hat a
  -
  \ee^{\ii\Omega t}\adag.
\end{equation}
Multiplicando ambas partes,
\begin{align}
  \hat H_{\vb r\cdot\vb E}^{(I)}(t)
  &=
  \hbar g
  \left(
    \ee^{\ii\omega t}\splus
    -
    \ee^{-\ii\omega t}\sminus
  \right)
  \left(
    \ee^{-\ii\Omega t}\hat a
    -
    \ee^{\ii\Omega t}\adag
  \right)                                                                    \\
  &=
  \hbar g
  \left[
    \ee^{-\ii(\Omega-\omega)t}\splus\hat a
    -
    \ee^{\ii(\Omega+\omega)t}\splus\adag
    -
    \ee^{-\ii(\Omega+\omega)t}\sminus\hat a
    +
    \ee^{\ii(\Omega-\omega)t}\sminus\adag
  \right].
\end{align}

\section{Aproximaci\'on de onda rotante}

Se define el desafine
\begin{equation}
  \Delta
  =
  \Omega-\omega.
\end{equation}
Los t\'erminos
\(\splus\adag\) y \(\sminus\hat a\) oscilan con frecuencia
\(\Omega+\omega\), que es grande cerca de resonancia. En integrales temporales
su contribuci\'on se promedia porque la fase cambia muy r\'apido. Los t\'erminos
\(\splus\hat a\) y \(\sminus\adag\) oscilan con \(\Delta\), que puede ser
peque\~no. Por eso se conservan solo estos t\'erminos lentamente oscilantes:
\begin{equation}
  \boxed{
  \hat H_{\mathrm{int}}^{(I)}(t)
  \simeq
  \hbar g
  \left(
    \ee^{-\ii\Delta t}\splus\hat a
    +
    \ee^{\ii\Delta t}\sminus\adag
  \right).
  }
\end{equation}
Esta es la aproximaci\'on de onda rotante. F\'isicamente conserva absorci\'on
\(\splus\hat a\), donde el \'atomo gana una excitaci\'on y el campo pierde un
fot\'on, y emisi\'on \(\sminus\adag\), donde el \'atomo pierde una excitaci\'on
y el campo gana un fot\'on.

Para regresar al cuadro de Schr\"odinger se usa la transformaci\'on inversa
\begin{equation}
  \hat H_{\vb r\cdot\vb E}^{(I)}
  =
  \ee^{\frac{\ii}{\hbar}\hat H_0t}
  \hat H_{\vb r\cdot\vb E}
  \ee^{-\frac{\ii}{\hbar}\hat H_0t},
  \qquad
  \hat H_{\vb r\cdot\vb E}
  =
  \ee^{-\frac{\ii}{\hbar}\hat H_0t}
  \hat H_{\vb r\cdot\vb E}^{(I)}
  \ee^{\frac{\ii}{\hbar}\hat H_0t}.
\end{equation}
Al hacerlo, se obtiene la interacci\'on resonante en el cuadro de
Schr\"odinger:
\begin{equation}
  \hat H_{\vb r\cdot\vb E}
  =
  \hbar g
  \left(
    \sminus\adag
    +
    \splus\hat a
  \right).
\end{equation}
Finalmente, el hamiltoniano para la interacci\'on de un \'atomo de dos niveles
con un modo del campo de radiaci\'on es
\begin{equation}
  \boxed{
  \hat H
  =
  \frac{\hat{\vb P}^{\,2}}{2M}
  +
  \hbar\Omega\,\adag\hat a
  +
  \frac{\hbar\omega}{2}\hat\sigma_z
  +
  \hbar g(\hat{\vb R})
  \left(
    \sminus\adag
    +
    \splus\hat a
  \right).
  }
\end{equation}
Si se desprecia el movimiento del centro de masa y se toma \(g\) constante, se
recupera el hamiltoniano de Jaynes--Cummings--Paul. La estructura f\'isica queda
clara: el campo y el \'atomo intercambian una excitaci\'on a la vez, y el
acoplamiento efectivo depende de la amplitud del campo de un fot\'on y del
elemento de matriz dipolar de la transici\'on.

\end{document}
